\section{Trigger (Dasar)}

\textbf{Trigger} adalah prosedur yang secara otomatis dijalankan (\textit{fired}) oleh MySQL ketika terjadi event tertentu pada tabel: \code{BEFORE INSERT}, \code{AFTER INSERT}, \code{BEFORE UPDATE}, \code{AFTER UPDATE}, \code{BEFORE DELETE}, atau \code{AFTER DELETE}. Trigger berguna untuk audit log (mencatat perubahan), validasi data di level database, atau memperbarui tabel lain secara otomatis. Dalam trigger, \code{NEW.kolom} merujuk ke nilai baru (INSERT/UPDATE) dan \code{OLD.kolom} merujuk ke nilai lama (UPDATE/DELETE) \cite{mysqlDocs}.

\begin{webcode}[language=SQL, caption={Trigger untuk audit log}]
-- Tabel audit log
CREATE TABLE produk_log (
  id        INT AUTO_INCREMENT PRIMARY KEY,
  produk_id INT, aksi VARCHAR(10),
  detail    TEXT, waktu DATETIME DEFAULT CURRENT_TIMESTAMP
);

-- Trigger: catat setiap UPDATE pada tabel produk
DELIMITER //
CREATE TRIGGER trg_produk_update
AFTER UPDATE ON produk
FOR EACH ROW
BEGIN
  INSERT INTO produk_log (produk_id, aksi, detail) VALUES (
    OLD.id, 'UPDATE',
    CONCAT('Harga: ', OLD.harga, ' -> ', NEW.harga)
  );
END //
DELIMITER ;

-- Ketika produk diupdate, log otomatis tercatat
UPDATE produk SET harga = 9000000 WHERE id = 1;
SELECT * FROM produk_log;
\end{webcode}

